\chapter{Red y seguridad}\label{cap:seguridad}

La seguridad está en tres capas, cada una contra un riesgo distinto. En una
red local de confianza alcanza con la primera, que no tiene costo; las otras
permiten exponer el mismo dispositivo en una red pública sin cambiar el
protocolo.

\begin{center}
\begin{tikzpicture}[node distance=4mm]
  \node[blk, minimum width=44mm, fill=m2teal!8] (t) {\textbf{Token}\\\scriptsize ¿quién sos?};
  \node[blk, minimum width=44mm, right=of t, fill=m2blue!8] (s) {\textbf{TLS}\\\scriptsize ¿alguien escucha?};
  \node[blk, minimum width=44mm, right=of s, fill=m2amber!12] (h) {\textbf{Defensas}\\\scriptsize ¿te pueden voltear?};
  \node[lbl, below=1mm of t] {siempre · costo nulo};
  \node[lbl, below=1mm of s] {opcional · $\sim$1 core};
  \node[lbl, below=1mm of h] {siempre · costo nulo};
\end{tikzpicture}
\end{center}

\section{Token de acceso}

Todo pedido, salvo \ep{GET}{/v1/health}, debe llevar el token:

\begin{lstlisting}[language=shell]
curl -H "Authorization: Bearer $(cat token)" http://192.168.18.5:8080/v1/device
\end{lstlisting}

\begin{itemize}
  \item Es un secreto aleatorio de 256 bits, generado en la primera
    instalación, en \cmd{/etc/mpeg2fpgad/token} (sólo legible por el servicio
    y root).
  \item Se compara en tiempo constante, para no revelar información a quien
    mida cuánto tarda la respuesta.
  \item Para rotarlo:
    \cmd{ssh root@placa python3 -m mpeg2fpgad -{}-rotate-token}, y reiniciar
    el servicio.
\end{itemize}

\begin{advertencia}
Sin TLS el token viaja sin cifrar: quien pueda observar el tráfico de la red
puede copiarlo. En una red de confianza es aceptable; en una red pública, no.
\end{advertencia}

\section{TLS}

Con \cmd{-{}-tls} el servicio atiende HTTPS en el puerto 8443 con un
certificado autofirmado generado en la placa. El cliente no confía en una
autoridad certificante sino en la \emph{huella} de ese certificado, obtenida
por un canal seguro:

\begin{lstlisting}[language=shell]
ssh root@placa python3 -m mpeg2fpgad --show-fingerprint
sha256:3f9a...
\end{lstlisting}

\begin{lstlisting}[language=Python]
dev = Device("192.168.18.5", token=token, tls_fingerprint="sha256:3f9a...")
\end{lstlisting}

Si la huella no coincide, la librería se niega a conectarse
(\cmd{FingerprintMismatch}): eso es lo que impide que un tercero se haga pasar
por el dispositivo.

El costo de TLS está medido en la sección de especificaciones. Para un
despliegue público con más caudal, la alternativa es terminar TLS en un equipo
delante del dispositivo (un \emph{reverse proxy}) y dejar la placa en una red
privada.

\section{Defensas básicas}

\begin{itemize}
  \item El servicio corre como usuario sin privilegios, con el sistema de
    archivos en sólo lectura, y recibe acceso sólo a los archivos del
    decodificador que usa.
  \item Diez tokens incorrectos seguidos desde una misma dirección la bloquean
    60 segundos (respuesta 429).
  \item A lo sumo 8 conexiones simultáneas; una conexión que no envía nada en
    30 segundos se cierra.
  \item Los endpoints de diagnóstico están apagados salvo \cmd{-{}-debug}.
  \item El clip sólo llega al decodificador: nada de lo que envía el cliente
    se usa como archivo o comando. Un clip malformado puede decodificarse mal,
    pero no afecta al dispositivo.
\end{itemize}
